\documentclass{article}
\usepackage[T1]{fontenc}
\usepackage{lmodern}
\usepackage{graphicx}
\usepackage{amsmath,amssymb}
\usepackage[a4paper,margin=2cm]{geometry}
\usepackage[absolute,overlay]{textpos}
\setlength{\TPHorizModule}{1mm}
\setlength{\TPVertModule}{1mm}
\usepackage[indonesian]{babel}
\usepackage{hyperref}
\setlength{\emergencystretch}{2em}
% unit: Halaman 73

\begin{document}

% === Halaman 73 (llm) ===

1. Program Steganografi dengan Python

\begin{itemize}
    \item Kode program Python berikut ini menyisipkan pesan dengan metode LSB.
    \item \textbf{A. Program penyisipan pesan}
    \begin{itemize}
        \item Input:
        \begin{enumerate}
            \item Citra berwarna (RGB)
            \item Pesan teks (diketik langsung dari keyboard, atau copas dari file teks)
        \end{enumerate}
        \item Output: citra stego (\textit{stego image})
    \end{itemize}
    \item \textbf{B. Program ekstraksi pesan}
    \begin{itemize}
        \item Input: Citra stego
        \item Output: pesan
    \end{itemize}
\end{itemize}

\end{document}
